\section{How to Rally}

\begin{rulebox}{General Rule}
The procedure by which disrupted units regain full effectiveness is called rallying. The Royalist player must make a rally check for all eligible disrupted Royalist units during the Royalist rally phase. The Allied player must make a rally check for all eligible disrupted Allied units during the Allied rally phase. No rally checks are possible at any other time during the game turn. Eligibility for rallying is defined in case 6.2.
\end{rulebox}

\ruleslabel{Procedure}

The active player executes the following steps in order for each of his eligible disrupted units:

\begin{steps}
  \item Find the morale rating of the unit for which the rally check is being made.
  \item Modify the printed morale rating by adding to it the command rating of any one leader of the active player's choice that is of the same force (color) as the unit and that is either stacked in the same hex with or adjacent to the unit.
  \item Add 2 to the number found in step 2 if the unit is a foot unit in a close hex, or add 1 if the unit (regardless of type) is in its own train hex. Subtract 2 if the unit (regardless of type) is in the enemy train hex. These modifiers are cumulative. The Royalist Train hex is hex 2810. The Allied Train hex is hex 0814.
  \item Roll one die, add 1 to the result if visibility is obscured or minimal, and compare the result to the modified morale rating found in step 3. If the die-roll result is less than or equal to the modified morale rating, the unit rallies; flip it to its ordered (front) side. If the result is greater, the unit remains disrupted. Another rally check must be made for the unit (if it is still eligible) during the next friendly rally phase. Die rolls greater than 6 are treated as 6.
\end{steps}

\ruleslabel{Cases}

\caseheading{6.1} The active player makes rally checks for his eligible disrupted units in any order he chooses.

However, he must make a rally check for each of his eligible disrupted units during each of his rally phases, and he cannot make a rally check for any ineligible disrupted unit.

\caseheading{6.2} A unit's eligibility for a rally check is affected by the proximity of enemy units and the morale of the army of which it is a part.

\begin{tightitems}
  \item Units are ineligible for a rally check when next to ordered enemy units, unless the disrupted unit is a foot unit in a woods, close, or train hex, and all the ordered enemy units are horse units.
  \item Units that are part of a demoralized army (13.0) are ineligible for a rally check, unless any of the following applies:
    \begin{enumerate}[label=\Alph*.,leftmargin=1.4em,itemsep=.15em,topsep=.15em]
      \item The unit is stacked in the same hex with a leader of the same force (color).
      \item The unit is one of the six horse units of Manchester's force.
      \item The unit is one of the seven foot units of Newcastle's force.
    \end{enumerate}
\end{tightitems}

Except as noted herein, all disrupted units are always eligible for a rally check during a friendly rally phase.
